\chapter{高频编码面试题}

\section{手写 Scaled Dot-Product Attention}
\begin{interview}
给定 \code{q}、\code{k}、\code{v}，手写缩放点积注意力。要求支持 mask，mask 中 \code{False} 的位置不能被关注。
\end{interview}
\begin{answer}
公式为
\[
  \operatorname{Attention}(Q,K,V)=\operatorname{softmax}\left(\frac{QK^\top}{\sqrt{d_k}}\right)V.
\]
\begin{lstlisting}
import math
import torch


def scaled_dot_product_attention(q, k, v, mask=None, dropout=None):
    """
    q: (..., tgt_len, d_k)
    k: (..., src_len, d_k)
    v: (..., src_len, d_v)
    mask: broadcastable to (..., tgt_len, src_len), True means keep
    """
    scores = torch.matmul(q, k.transpose(-2, -1))
    scores = scores / math.sqrt(q.size(-1))
    if mask is not None:
        scores = scores.masked_fill(~mask, torch.finfo(scores.dtype).min)
    attn = torch.softmax(scores, dim=-1)
    if dropout is not None:
        attn = dropout(attn)
    out = torch.matmul(attn, v)
    return out, attn
\end{lstlisting}
\paragraph{要点} \code{k.transpose(-2, -1)} 只交换最后两维；mask 要能广播到注意力分数形状；softmax 沿 key 维，即 \code{dim=-1}。
\end{answer}
\begin{trap}
不要先 softmax 再 mask；被屏蔽位置应在 softmax 前填成足够小的值。
\end{trap}

\section{手写 Multi-Head Attention}
\begin{interview}
用 \code{nn.Linear} 实现多头注意力：线性投影、拆头、缩放点积注意力、合头、输出投影。
\end{interview}
\begin{answer}
多头注意力把 \code{embed\_dim} 分成多个 \code{head\_dim}，每个头独立做 attention。
\begin{lstlisting}
import math
import torch
from torch import nn


class MultiHeadAttention(nn.Module):
    def __init__(self, embed_dim, num_heads, dropout=0.0):
        super().__init__()
        if embed_dim % num_heads != 0:
            raise ValueError("embed_dim must be divisible by num_heads")
        self.embed_dim = embed_dim
        self.num_heads = num_heads
        self.head_dim = embed_dim // num_heads
        self.q_proj = nn.Linear(embed_dim, embed_dim)
        self.k_proj = nn.Linear(embed_dim, embed_dim)
        self.v_proj = nn.Linear(embed_dim, embed_dim)
        self.out_proj = nn.Linear(embed_dim, embed_dim)
        self.dropout = nn.Dropout(dropout)

    def split_heads(self, x):
        bsz, seq_len, _ = x.shape
        x = x.view(bsz, seq_len, self.num_heads, self.head_dim)
        return x.transpose(1, 2)

    def merge_heads(self, x):
        bsz, heads, seq_len, head_dim = x.shape
        x = x.transpose(1, 2).contiguous()
        return x.view(bsz, seq_len, heads * head_dim)

    def forward(self, query, key, value, mask=None):
        q = self.split_heads(self.q_proj(query))
        k = self.split_heads(self.k_proj(key))
        v = self.split_heads(self.v_proj(value))
        scores = q @ k.transpose(-2, -1) / math.sqrt(self.head_dim)
        if mask is not None:
            if mask.dim() == 3:
                mask = mask.unsqueeze(1)
            scores = scores.masked_fill(~mask, torch.finfo(scores.dtype).min)
        attn = self.dropout(torch.softmax(scores, dim=-1))
        out = self.merge_heads(attn @ v)
        return self.out_proj(out), attn
\end{lstlisting}
\paragraph{要点} 输入常用 \code{(batch, seq\_len, embed\_dim)}；拆头后是 \code{(batch, heads, seq\_len, head\_dim)}；\code{transpose} 后再 \code{view} 前通常要 \code{contiguous()}。
\end{answer}

\section{用 \code{nn.Parameter} 手写 Linear 层}
\begin{interview}
不用 \code{nn.Linear}，只用 \code{nn.Parameter} 手写全连接层，包含初始化和 \code{forward}。
\end{interview}
\begin{answer}
线性层计算 $Y=XW^\top+b$。PyTorch 线性层权重形状为 \code{(out\_features, in\_features)}。
\begin{lstlisting}
import math
import torch
from torch import nn


class MyLinear(nn.Module):
    def __init__(self, in_features, out_features, bias=True):
        super().__init__()
        self.in_features = in_features
        self.out_features = out_features
        self.weight = nn.Parameter(torch.empty(out_features, in_features))
        if bias:
            self.bias = nn.Parameter(torch.empty(out_features))
        else:
            self.register_parameter("bias", None)
        self.reset_parameters()

    def reset_parameters(self):
        bound = 1.0 / math.sqrt(self.in_features)
        nn.init.uniform_(self.weight, -bound, bound)
        if self.bias is not None:
            nn.init.uniform_(self.bias, -bound, bound)

    def forward(self, x):
        y = x.matmul(self.weight.t())
        if self.bias is not None:
            y = y + self.bias
        return y
\end{lstlisting}
\paragraph{要点} 可训练张量必须是 \code{nn.Parameter}；无 bias 时用 \code{register\_parameter("bias", None)}；bias 自动广播到最后一维。
\end{answer}

\section{手写 BatchNorm1d}
\begin{interview}
手写 \code{BatchNorm1d}，区分训练/推理分支，训练时更新 \code{running\_mean} 与 \code{running\_var}。
\end{interview}
\begin{answer}
BatchNorm 先归一化再仿射变换：$\hat{x}=(x-\mu)/\sqrt{\sigma^2+\epsilon}$，$y=\gamma\hat{x}+\beta$。
\begin{lstlisting}
import torch
from torch import nn


class MyBatchNorm1d(nn.Module):
    def __init__(self, num_features, eps=1e-5, momentum=0.1, affine=True):
        super().__init__()
        self.eps = eps
        self.momentum = momentum
        self.affine = affine
        if affine:
            self.weight = nn.Parameter(torch.ones(num_features))
            self.bias = nn.Parameter(torch.zeros(num_features))
        else:
            self.register_parameter("weight", None)
            self.register_parameter("bias", None)
        self.register_buffer("running_mean", torch.zeros(num_features))
        self.register_buffer("running_var", torch.ones(num_features))

    def forward(self, x):
        if x.dim() not in (2, 3):
            raise ValueError("expected 2D or 3D input")
        reduce_dims = (0,) if x.dim() == 2 else (0, 2)
        shape = (1, -1) if x.dim() == 2 else (1, -1, 1)
        if self.training:
            mean = x.mean(dim=reduce_dims, keepdim=True)
            var = x.var(dim=reduce_dims, unbiased=False, keepdim=True)
            with torch.no_grad():
                self.running_mean.mul_(1 - self.momentum).add_(
                    self.momentum * mean.view(-1)
                )
                self.running_var.mul_(1 - self.momentum).add_(
                    self.momentum * var.view(-1)
                )
        else:
            mean = self.running_mean.view(shape)
            var = self.running_var.view(shape)
        y = (x - mean) / torch.sqrt(var + self.eps)
        if self.affine:
            y = y * self.weight.view(shape) + self.bias.view(shape)
        return y
\end{lstlisting}
\paragraph{要点} running stats 是 buffer 不是参数；\code{(N,C,L)} 要在 \code{N} 和 \code{L} 维统计；更新 buffer 放进 \code{torch.no\_grad()}。
\end{answer}
\begin{trap}
官方实现对 running variance 的无偏估计更细；面试重点通常是分支、buffer 与广播。
\end{trap}

\section{手写 Dropout}
\begin{interview}
实现 inverted dropout：训练时随机置零并缩放，推理时直接返回输入。
\end{interview}
\begin{answer}
inverted dropout 在训练阶段除以 \code{1-p}，推理阶段无需缩放。
\begin{lstlisting}
import torch
from torch import nn


class MyDropout(nn.Module):
    def __init__(self, p=0.5):
        super().__init__()
        if p < 0.0 or p >= 1.0:
            raise ValueError("p must be in [0, 1)")
        self.p = p

    def forward(self, x):
        if not self.training or self.p == 0.0:
            return x
        keep_prob = 1.0 - self.p
        mask = torch.rand_like(x) < keep_prob
        return x * mask.to(x.dtype) / keep_prob
\end{lstlisting}
\paragraph{要点} 训练时输出期望保持不变；推理时不要再乘 \code{1-p}；用 \code{torch.rand\_like(x)} 可自动匹配 device。
\end{answer}

\section{手写交叉熵损失}
\begin{interview}
手写多分类交叉熵损失，要求用 \code{log\_softmax + nll}，输入是 logits 和类别索引。
\end{interview}
\begin{answer}
不要先 softmax 再取 log；稳定写法是直接计算 \code{log\_softmax}。
\begin{lstlisting}
import torch
import torch.nn.functional as F


def cross_entropy_from_logits(logits, target, reduction="mean"):
    log_probs = F.log_softmax(logits, dim=-1)
    nll = -log_probs.gather(1, target.view(-1, 1)).squeeze(1)
    if reduction == "mean":
        return nll.mean()
    if reduction == "sum":
        return nll.sum()
    if reduction == "none":
        return nll
    raise ValueError("invalid reduction")
\end{lstlisting}
\paragraph{要点} logits 不需要预先归一化；target 形状通常是 \code{(batch,)} 且 dtype 为 \code{torch.long}；\code{gather} 的 index 维度要匹配。
\end{answer}

\section{手写完整训练循环}
\begin{interview}
实现 \code{train\_one\_epoch(model, loader, opt, loss\_fn, device)}，包含训练模式、搬运数据、前向、反向、更新和平均损失统计。
\end{interview}
\begin{answer}
训练循环的标准顺序是清梯度、前向、算损失、反向、更新参数。
\begin{lstlisting}
import torch


def train_one_epoch(model, loader, opt, loss_fn, device):
    model.train()
    total_loss = 0.0
    total_count = 0
    for batch in loader:
        if isinstance(batch, (tuple, list)):
            x, y = batch[0], batch[1]
        else:
            x, y = batch["x"], batch["y"]
        x = x.to(device)
        y = y.to(device)
        opt.zero_grad(set_to_none=True)
        pred = model(x)
        loss = loss_fn(pred, y)
        loss.backward()
        opt.step()
        batch_size = x.size(0)
        total_loss += loss.detach().item() * batch_size
        total_count += batch_size
    return total_loss / max(total_count, 1)
\end{lstlisting}
\paragraph{要点} 顺序是 \code{zero\_grad -> forward -> loss -> backward -> step}；统计 loss 时用 \code{detach().item()}；若 loss 是 batch mean，累计时乘 batch size。
\end{answer}
\begin{trap}
评估循环才使用 \code{model.eval()} 和 \code{torch.no\_grad()}；训练循环不能包 \code{no\_grad}。
\end{trap}

\section{张量操作题：不用 Python for 循环}
\begin{interview}
不用 Python for 循环实现：按行 softmax、按标签取概率、广播计算成对欧氏距离、topk、masked\_fill。
\end{interview}
\begin{answer}
这类题考维度意识，应优先用广播、\code{gather}、\code{topk}、\code{masked\_fill} 等原语。
\begin{lstlisting}
import torch
import torch.nn.functional as F


def row_softmax(x):
    return torch.softmax(x, dim=1)


def take_class_prob(probs, labels):
    return probs.gather(1, labels.view(-1, 1)).squeeze(1)


def pairwise_l2_distance(x, y):
    diff = x[:, None, :] - y[None, :, :]
    return torch.sqrt((diff * diff).sum(dim=-1))


def pairwise_l2_distance_fast(x, y):
    x2 = (x * x).sum(dim=1, keepdim=True)
    y2 = (y * y).sum(dim=1).unsqueeze(0)
    dist2 = (x2 + y2 - 2.0 * x @ y.t()).clamp_min(0.0)
    return torch.sqrt(dist2)


def topk_values_indices(x, k):
    return torch.topk(x, k=k, dim=1)


def mask_logits(logits, valid_mask):
    return logits.masked_fill(~valid_mask, torch.finfo(logits.dtype).min)


def demo_tensor_ops():
    logits = torch.tensor([[1.0, 2.0, 3.0], [3.0, 1.0, 0.0]])
    labels = torch.tensor([2, 0])
    probs = F.softmax(logits, dim=1)
    chosen = take_class_prob(probs, labels)
    x = torch.randn(4, 5)
    y = torch.randn(6, 5)
    dist = pairwise_l2_distance_fast(x, y)
    values, indices = topk_values_indices(logits, k=2)
    mask = torch.tensor([[True, True, False], [True, False, True]])
    masked = mask_logits(logits, mask)
    return chosen, dist, values, indices, masked
\end{lstlisting}
\paragraph{要点} \code{gather} 适合按样本取类别；\code{x[:, None, :]} 和 \code{y[None, :, :]} 可广播；\code{topk} 返回 values 和 indices；先说明 mask 中 \code{True} 的含义。
\end{answer}

\section{可选：Label Smoothing 交叉熵}
\begin{interview}
实现 label smoothing 版本交叉熵。输入 logits 和类别索引 target，平滑系数为 \code{smoothing}。
\end{interview}
\begin{answer}
label smoothing 把少量概率质量分给非正确类别，可缓解模型过度自信。
\begin{lstlisting}
import torch
import torch.nn.functional as F

def label_smoothing_cross_entropy(logits, target, smoothing=0.1):
    num_classes = logits.size(-1)
    log_probs = F.log_softmax(logits, dim=-1)
    with torch.no_grad():
        true_dist = torch.full_like(log_probs, smoothing / (num_classes - 1))
        true_dist.scatter_(1, target.view(-1, 1), 1.0 - smoothing)
    loss = -(true_dist * log_probs).sum(dim=-1)
    return loss.mean()
\end{lstlisting}
\paragraph{要点} \code{scatter\_} 写入正确类别概率；构造软标签不需要梯度；实际任务应保证 \code{num\_classes > 1}。
\end{answer}
\begin{quickcard}
\textbf{手写模板速记}
\begin{itemize}
  \item Attention/MHA：$QK^\top/\sqrt{d_k}$，mask 后对最后一维 softmax；多头要线性投影、拆头、合头。
  \item Linear/BatchNorm/Dropout：参数用 \code{nn.Parameter}，running stats 用 buffer，inverted dropout 训练除以 \code{1-p}。
  \item CrossEntropy/Train loop：稳定写法是 \code{log\_softmax + gather}；训练顺序是清梯度、前向、反向、更新。
  \item Tensor ops：多用广播、\code{gather}、\code{topk}、\code{masked\_fill}，并先说明 mask 中 \code{True} 的语义。
\end{itemize}
\end{quickcard}
